(1) 两条中子干涉路径在水平长度 \(L\) 上分别处于高度 \(0\) 和 \(H\)，竖直连接段贡献相同，重力加速度 \(g\)，中子质量 \(m\)、速率 \(v\)。从作用量差求相位差。
(2) 用德布罗意波长 \(\lambda=h/(mv)\) 与面积 \(A=HL\) 表示相位。
(3) 求将干涉仪绕水平轴倾斜角 \(\alpha\) 时的相位变化和条纹周期。
(4) 讨论 \(g\to0\)、\(A\to0\)、高速极限及连接段抵消假设。
